%----------------------------------------------------------------------
\subsection{Calculation of the single-quark incoming state}
\label{sec:wf-q}
%----------------------------------------------------------------------

The pure Yang--Mills construction has no analogue of this sector. A soft quark
$\ket{\qk{\beta'}{\lambda'}{l'}}$ emits and absorbs soft gluons and rescatters on the valence
charge; these processes fix the coefficients $\mf F$, $\mf P$, $\mf Q_{1}$ and $\mf K_{1}$. The
results for an incoming antiquark are given at the end of this subsection. The outgoing quark is
$\qk{\beta}{\lambda}{l}$.

\subsubsection{Quark and one gluon outgoing}

To order $g$ the quark emits a soft gluon.
\begin{center}
\begin{tikzpicture}
\draw[qk] (0,1.0) -- (1.7,1.0); \node[left] at (0,1.0) {\scriptsize$\qk{\beta'}{\lambda'}{l'}$};
\draw[qk] (1.7,1.0) -- (3.6,1.6); \node[right] at (3.6,1.6) {\scriptsize$\qk{\beta}{\lambda}{l}$};
\draw[gl] (1.7,1.0) -- (3.6,0.3); \node[right] at (3.6,0.3) {\scriptsize$\gl{b}{j}{p}$};
\draw[cut] (2.6,0.0) -- (2.6,1.95);
\end{tikzpicture}
\end{center}
\begin{multline}
\ket{(\Psi^{\beta'}_{\lambda'}(l'))_{qg}}=\int\!d^{3}l\,d^{3}p\;\ket{\qk{\beta}{\lambda}{l}\gl{b}{j}{p}}\,
\frac{\bra{\qk{\beta}{\lambda}{l}\gl{b}{j}{p}}H_{gq\bar q}\ket{\qk{\beta'}{\lambda'}{l'}}}{\eps_{l'}-\eps_{l}-\eps_{p}}\\
=\frac{g}{8\pi^{3/2}}\int\!d^{3}l\,d^{3}p\;\ket{qg}\,
\frac{t^{b}_{\beta\beta'}\,\dthree(l'-l-p)\,\Ph^{j}_{\lambda\lambda'}(p;l,l')}{\sqrt{p^{+}}\,\big(\eps_{l'}-\eps_{l}-\eps_{p}\big)} ,
\label{eq:Pq-qg}
\end{multline}
with the matrix element \eqref{eq:MEqg}. The Weizs\"acker--Williams gluon emitted by $\rho$ with
the quark as a spectator is the disconnected $A$ term of \eqref{eq:Omq}.

\subsubsection{One quark outgoing}

At order $g^{2}$ the soft quark exchanges a gluon with the valence charge. Three diagrams
contribute to the transition $l'\to l$ with momentum transfer to the valence; we write
\begin{equation}
\Delta\equiv E_{i}-E_{f}=\eps_{l'}-\eps_{l}.
\end{equation}

\paragraph{(a) The quark emits a gluon which is absorbed by the valence.}
\begin{center}
\begin{tikzpicture}
\draw[val] (0,0) -- (3.9,0); \node[left] at (-0.05,0) {\scriptsize$\rho$};
\draw[qk] (0,1.6) -- (1.3,1.6); \node[left] at (0,1.6) {\scriptsize$\qk{\beta'}{\lambda'}{l'}$};
\draw[qk] (1.3,1.6) -- (3.9,1.6); \node[right] at (3.9,1.6) {\scriptsize$\qk{\beta}{\lambda}{l}$};
\draw[gl] (1.3,1.6) -- (2.8,0); \node[right] at (2.15,0.75) {\scriptsize$\gl{b}{j}{p}$};
\draw[cut] (2.05,-0.3) -- (2.05,1.95); \draw[cut] (3.3,-0.3) -- (3.3,1.95);
\end{tikzpicture}
\end{center}
Here $p=l'-l$ and the process requires $l'^{+}>l^{+}$:
\begin{equation}
\ket{(\Psi^{\beta'}_{\lambda'}(l'))^{1}_{q\,\rho}}=\int\!d^{3}l\,d^{3}p\;\ket{\qk{\beta}{\lambda}{l}}\;
\frac{\bra{\qk{\beta}{\lambda}{l}}H_{g}\ket{\qk{\beta}{\lambda}{l}\gl{b}{j}{p}}}{\Delta}\;
\frac{\bra{\qk{\beta}{\lambda}{l}\gl{b}{j}{p}}H_{gq\bar q}\ket{\qk{\beta'}{\lambda'}{l'}}}{\Delta-\eps_{p}} .
\end{equation}
Using \eqref{eq:MEg2} and \eqref{eq:MEqg},
\begin{equation}
\ket{(\Psi^{\beta'}_{\lambda'}(l'))^{1}_{q\,\rho}}=\frac{g^{2}}{32\pi^{3}}\int\!d^{3}l\;\ket{\qk{\beta}{\lambda}{l}}\;
\frac{t^{b}_{\beta\beta'}\,\rho^{b}(p)\,p^{j}\Ph^{j}_{\lambda\lambda'}(p;l,l')}{(p^{+})^{2}\,(\Delta-\eps_{p})\,\Delta}\,
\theta(l'^{+}-l^{+}) .
\label{eq:P1a}
\end{equation}

\paragraph{(b) The valence emits a gluon which is absorbed by the quark.}
\begin{center}
\begin{tikzpicture}
\draw[val] (0,0) -- (3.9,0); \node[left] at (-0.05,0) {\scriptsize$\rho$};
\draw[qk] (0,1.6) -- (2.6,1.6); \node[left] at (0,1.6) {\scriptsize$\qk{\beta'}{\lambda'}{l'}$};
\draw[qk] (2.6,1.6) -- (3.9,1.6); \node[right] at (3.9,1.6) {\scriptsize$\qk{\beta}{\lambda}{l}$};
\draw[gl] (1.1,0) -- (2.6,1.6); \node[left] at (1.75,0.85) {\scriptsize$\gl{b}{j}{p}$};
\draw[cut] (1.85,-0.3) -- (1.85,1.95); \draw[cut] (3.3,-0.3) -- (3.3,1.95);
\end{tikzpicture}
\end{center}
Now $p=l-l'$, the process requires $l^{+}>l'^{+}$ and the intermediate state
$\ket{\qk{\beta'}{\lambda'}{l'}\gl{b}{j}{p}}$ has $E_{i}-E_{m}=-\eps_{p}$:
\begin{equation}
\ket{(\Psi^{\beta'}_{\lambda'}(l'))^{2}_{q\,\rho}}=-\frac{g^{2}}{32\pi^{3}}\int\!d^{3}l\;\ket{\qk{\beta}{\lambda}{l}}\;
\frac{t^{b}_{\beta\beta'}\,\rho^{b}(-p)\,p^{j}\Ph^{j}_{\lambda\lambda'}(p;l,l')}{(p^{+})^{2}\,\eps_{p}\,\Delta}\,
\theta(l^{+}-l'^{+}) ,
\label{eq:P1b}
\end{equation}
where the absorption vertex is $\big[\Ph^{j}_{\lambda'\lambda}(p;l',l)\big]^{*}=\Ph^{j}_{\lambda\lambda'}(p;l,l')$.

\paragraph{(c) Instantaneous gluon exchange with the valence.}
\begin{center}
\begin{tikzpicture}
\draw[val] (0,0) -- (3.9,0); \node[left] at (-0.05,0) {\scriptsize$\rho$};
\draw[qk] (0,1.4) -- (3.9,1.4); \node[left] at (0,1.4) {\scriptsize$\qk{\beta'}{\lambda'}{l'}$};
\node[right] at (3.9,1.4) {\scriptsize$\qk{\beta}{\lambda}{l}$};
\draw[ins] (1.9,0) -- (1.9,1.4); \ibar{1.9}{0.7}
\draw[cut] (3.2,-0.3) -- (3.2,1.75);
\end{tikzpicture}
\end{center}
This is the eikonal colour rotation of the soft quark by the $\rho\cdot J_{q}$ term of
\eqref{eq:Hqinst}; with \eqref{eq:MEinstq},
\begin{multline}
\ket{(\Psi^{\beta'}_{\lambda'}(l'))^{3}_{q\,\rho}}=\int\!d^{3}l\;\ket{\qk{\beta}{\lambda}{l}}\,
\frac{\bra{\qk{\beta}{\lambda}{l}}H_{q-{\rm inst}}\ket{\qk{\beta'}{\lambda'}{l'}}}{\Delta}\\
=\frac{g^{2}}{8\pi^{3}}\int\!d^{3}l\;\ket{\qk{\beta}{\lambda}{l}}\,
\frac{t^{b}_{\beta\beta'}\,\rho^{b}(l'-l)\,\delta_{\lambda\lambda'}}{(l'^{+}-l^{+})^{2}\,\Delta} .
\label{eq:P1c}
\end{multline}

\paragraph{(d) Self-energy.} The quark may also emit and reabsorb a soft gluon. This contribution
is proportional to $\dthree(l-l')$, has $E_{i}=E_{f}$, and enters only the normalization,
discussed below.

\paragraph{Simplification.} For $p=l'-l$ the vertex satisfies
\begin{equation}
p^{j}\Gam^{j}(p;l,l')=4\eps_{p}-2\Delta-\frac{p^{+}}{l^{+}l'^{+}}(\bm\sigma\!\cdot\!\bm l)(\bm\sigma\!\cdot\!\bm l'),
\label{eq:pGam1}
\end{equation}
and for $p=l-l'$, $p^{j}\Gam^{j}(p;l,l')=4\eps_{p}+2\Delta-\tfrac{p^{+}}{l^{+}l'^{+}}(\bm\sigma\!\cdot\!\bm l)(\bm\sigma\!\cdot\!\bm l')$.
Writing $4\eps_{p}-2\Delta=-4(\Delta-\eps_{p})+2\Delta$ in \eqref{eq:P1a}, and
$4\eps_{p}+2\Delta$ unchanged in \eqref{eq:P1b}, the terms proportional to
$\delta_{\lambda\lambda'}/\Delta$ cancel the instantaneous contribution \eqref{eq:P1c} in both
kinematic regions, exactly as in the vacuum sector. The sum is
\begin{align}
\ket{(\Psi^{\beta'}_{\lambda'}(l'))_{q\,\rho}}=\;&\frac{g^{2}}{32\pi^{3}}\int\!d^{3}l\;\ket{\qk{\beta}{\lambda}{l}}\,
\frac{t^{b}_{\beta\beta'}\rho^{b}(l'-l)}{(p^{+})^{2}}\,\chi^{\dg}_{\lambda}\Bigg\{
\nonumber\\&\hspace{0.6cm}
\theta(l'^{+}-l^{+})\left[\frac{2}{\Delta-\eps_{p}}
-\frac{p^{+}\,(\bm\sigma\!\cdot\!\bm l)(\bm\sigma\!\cdot\!\bm l')}{l^{+}l'^{+}\,(\Delta-\eps_{p})\,\Delta}\right]
\nonumber\\
&\hspace{0.6cm}-\theta(l^{+}-l'^{+})\left[\frac{2}{\eps_{p}}
-\frac{p^{+}\,(\bm\sigma\!\cdot\!\bm l)(\bm\sigma\!\cdot\!\bm l')}{l^{+}l'^{+}\,\eps_{p}\,\Delta}\right]\Bigg\}\chi_{\lambda'},
\label{eq:Pq-q}
\end{align}
where $p=\pm(l'-l)$ is the (positive plus momentum) gluon exchanged in the two regions. The
overall denominator $\Delta=\eps_{l'}-\eps_{l}$ survives in the helicity dependent term: the
coefficient $\mf P$ has an on-shell pole, exactly as the coefficient $B_{3}$ of Paper~I, and its
residue is the amplitude for the scattering of a soft quark on the valence charge.

\subsubsection{Quark and two gluons outgoing}

At order $g^{2}$ the transition $q(l')\to q(l)\,\gl{c}{k}{q}\gl{d}{n}{r}$ receives five
contributions:
\begin{center}
\begin{tikzpicture}[scale=0.85]
\draw[qk] (0,1.2) -- (1.0,1.2); \draw[qk] (1.0,1.2) -- (2.0,1.5); \draw[qk] (2.0,1.5) -- (3.0,1.8);
\draw[gl] (1.0,1.2) -- (3.0,0.5); \draw[gl] (2.0,1.5) -- (3.0,1.15);
\node at (1.5,-0.3) {\scriptsize (a)};
\end{tikzpicture}\hspace{0.35cm}
\begin{tikzpicture}[scale=0.85]
\draw[qk] (0,1.2) -- (1.0,1.2); \draw[qk] (1.0,1.2) -- (3.0,1.8);
\draw[gl] (1.0,1.2) -- (2.0,0.8); \draw[gl] (2.0,0.8) -- (3.0,1.1); \draw[gl] (2.0,0.8) -- (3.0,0.3);
\node at (1.5,-0.3) {\scriptsize (b)};
\end{tikzpicture}\hspace{0.35cm}
\begin{tikzpicture}[scale=0.85]
\draw[val] (0,0) -- (3.0,0);
\draw[qk] (0,1.4) -- (1.4,1.4); \draw[qk] (1.4,1.4) -- (3.0,1.8);
\draw[gl] (1.4,1.4) -- (3.0,1.0); \draw[gl] (0.9,0) -- (3.0,0.6);
\node at (1.5,-0.45) {\scriptsize (c)};
\end{tikzpicture}\hspace{0.35cm}
\begin{tikzpicture}[scale=0.85]
\draw[qk] (0,1.3) -- (1.1,1.3); \draw[qk] (1.6,1.3) -- (3.0,1.8);
\draw[ins] (1.1,1.3) -- (1.6,1.3); \ibarh{1.35}{1.3}
\draw[gl] (1.1,1.3) -- (3.0,0.4); \draw[gl] (1.6,1.3) -- (3.0,1.05);
\node at (1.5,-0.3) {\scriptsize (d)};
\end{tikzpicture}\hspace{0.35cm}
\begin{tikzpicture}[scale=0.85]
\draw[qk] (0,1.4) -- (1.4,1.4); \draw[qk] (1.4,1.4) -- (3.0,1.8);
\draw[ins] (1.4,1.4) -- (1.4,0.7); \ibar{1.4}{1.05}
\draw[gl] (1.4,0.7) -- (3.0,1.0); \draw[gl] (1.4,0.7) -- (3.0,0.3);
\node at (1.5,-0.3) {\scriptsize (e)};
\end{tikzpicture}
\end{center}
(a) two sequential emissions by the quark, in either order; (b) emission of a gluon which splits
into two through $H_{ggg}$; (c) emission of one gluon by the quark and of the other by the
valence charge; (d) instantaneous quark exchange; (e) instantaneous gluon exchange between the
quark and gluon currents. With $D_{f}=\eps_{l'}-\eps_{l}-\eps_{q}-\eps_{r}$,
\begin{align}
\ket{(\Psi^{\beta'}_{\lambda'}(l'))_{qgg}}=\;&\frac12\int\!d^{3}l\,d^{3}q\,d^{3}r\;\ket{\qk{\beta}{\lambda}{l}\gl{c}{k}{q}\gl{d}{n}{r}}
\nonumber\\&\times\Bigg\{\int\!d^{3}l_{1}\,\frac{\bra{\qk{\beta}{\lambda}{l}\gl{c}{k}{q}}H_{gq\bar q}\ket{q(l_{1})}}{D_{f}}\,
\frac{\bra{q(l_{1})\gl{d}{n}{r}}H_{gq\bar q}\ket{\qk{\beta'}{\lambda'}{l'}}}{\eps_{l'}-\eps_{l_{1}}-\eps_{r}}
+(q\leftrightarrow r)
\nonumber\\
&+\int\!d^{3}s\,\frac{\bra{\gl{c}{k}{q}\gl{d}{n}{r}}H_{ggg}\ket{g(s)}}{D_{f}}\,
\frac{\bra{\qk{\beta}{\lambda}{l}g(s)}H_{gq\bar q}\ket{\qk{\beta'}{\lambda'}{l'}}}{\eps_{l'}-\eps_{l}-\eps_{s}}
\nonumber\\
&+\frac{\bra{q\,gg}\big(H_{ggq\bar q}+H_{q-{\rm inst}}\big)\ket{q}}{D_{f}}\Bigg\}
\nonumber\\
&+\ket{(\Psi^{\beta'}_{\lambda'}(l'))_{qgg\,\rho}} ,
\label{eq:Pq-qgg}
\end{align}
where the factor $\tfrac12$ compensates the double counting of the identical gluons, the
instantaneous matrix elements are \eqref{eq:MEq2g} and \eqref{eq:MEJJq2g}, and the $\rho$ dependent
component (c) factorizes by \eqref{eq:factor} into the product of \eqref{eq:Pq-qg} and the
Weizs\"acker--Williams amplitude. The diagrams (a)--(e) are related by crossing to the diagrams
(a)--(g) of the $q\bar q g$ component of the one gluon state.

\subsubsection{Quark and a quark--antiquark pair outgoing}

At order $g^{2}$ the quark can also emit a gluon which splits into a pair, or create a pair
through the instantaneous $J_{q}\cdot J_{q}$ interaction:
\begin{center}
\begin{tikzpicture}
\draw[qk] (0,1.6) -- (1.2,1.6); \node[left] at (0,1.6) {\scriptsize$\qk{\beta'}{\lambda'}{l'}$};
\draw[qk] (1.2,1.6) -- (3.9,2.0); \node[right] at (3.9,2.0) {\scriptsize$\qk{\beta}{\lambda}{l}$};
\draw[gl] (1.2,1.6) -- (2.4,0.9); \node[below left] at (1.85,1.2) {\scriptsize$k$};
\draw[qk] (2.4,0.9) -- (3.9,1.25); \node[right] at (3.9,1.25) {\scriptsize$\qk{\beta_{2}}{\lambda_{2}}{l_{2}}$};
\draw[qk] (3.9,0.35) -- (2.4,0.9); \node[right] at (3.9,0.35) {\scriptsize$\aqk{\alpha}{\lambda_{m}}{m}$};
\draw[cut] (1.8,0.1) -- (1.8,2.3); \draw[cut] (3.2,0.1) -- (3.2,2.3);
\end{tikzpicture}
\hspace{0.8cm}
\begin{tikzpicture}
\draw[qk] (0,1.6) -- (1.8,1.6); \node[left] at (0,1.6) {\scriptsize$\qk{\beta'}{\lambda'}{l'}$};
\draw[qk] (1.8,1.6) -- (3.9,2.0); \node[right] at (3.9,2.0) {\scriptsize$\qk{\beta}{\lambda}{l}$};
\draw[ins] (1.8,1.6) -- (1.8,0.8); \ibar{1.8}{1.2}
\draw[qk] (1.8,0.8) -- (3.9,1.2); \node[right] at (3.9,1.2) {\scriptsize$\qk{\beta_{2}}{\lambda_{2}}{l_{2}}$};
\draw[qk] (3.9,0.35) -- (1.8,0.8); \node[right] at (3.9,0.35) {\scriptsize$\aqk{\alpha}{\lambda_{m}}{m}$};
\draw[cut] (3.2,0.1) -- (3.2,2.3);
\end{tikzpicture}
\end{center}
\begin{align}
\ket{(\Psi^{\beta'}_{\lambda'}(l'))_{qq\bar q}}=\;&\frac12\int\!d^{3}l\,d^{3}l_{2}\,d^{3}m\;\ket{\qk{\beta}{\lambda}{l}\qk{\beta_{2}}{\lambda_{2}}{l_{2}}\aqk{\alpha}{\lambda_{m}}{m}}
\nonumber\\&\times\Bigg\{\int\!d^{3}k\,
\frac{\bra{q(l_{2})\bar q(m)}H_{gq\bar q}\ket{g(k)}\bra{q(l)g(k)}H_{gq\bar q}\ket{q(l')}}{D_{f}\,\big(\eps_{l'}-\eps_{l}-\eps_{k}\big)}
\nonumber\\
&+\frac{\bra{q(l)q(l_{2})\bar q(m)}H_{q-{\rm inst}}\ket{q(l')}}{D_{f}}\Bigg\}
\nonumber\\&-\big((\beta,\lambda,l)\leftrightarrow(\beta_{2},\lambda_{2},l_{2})\big),
\label{eq:Pq-qqq}
\end{align}
with $D_{f}=\eps_{l'}-\eps_{l}-\eps_{l_{2}}-\eps_{m}$ and the matrix element \eqref{eq:MEJJqqq}. The
antisymmetrization accounts for the two identical quarks in the final state. The
$\rho$ dependent pair creation off the valence charge with a spectator quark is the
disconnected $\mf A_{1}$ term.

\subsubsection{Single-antiquark incoming state}

The antiquark sector follows from the quark sector by the replacements dictated by the
antiquark vertex \eqref{eq:MEaqg}: the colour matrix is replaced by
$t^{b}_{\beta\beta'}\to-t^{b}_{\alpha'\alpha}$ (incoming colour index first), and the spinor
bilinear of the quark line by that of the antiquark line, in which the incoming helicity stands
on the left, $\Ph^{j}_{\lambda\lambda'}(p;l,l')\to\Ph^{j}_{\lambda'\lambda}(p;m',m)$. For an incoming
antiquark $\ket{\aqk{\alpha'}{\lambda'}{m'}}$ and outgoing $\aqk{\alpha}{\lambda}{m}$ one finds
\begin{align}
\ket{(\Psi^{\alpha'}_{\lambda'}(m'))_{\bar qg}}&=-\frac{g}{8\pi^{3/2}}\int\!d^{3}m\,d^{3}p\;\ket{\bar q g}\,
\frac{t^{b}_{\alpha'\alpha}\,\dthree(m'-m-p)\,\Ph^{j}_{\lambda'\lambda}(p;m',m)}{\sqrt{p^{+}}\,\big(\eps_{m'}-\eps_{m}-\eps_{p}\big)},
\label{eq:Pqb-qbg}\\
\ket{(\Psi^{\alpha'}_{\lambda'}(m'))_{\bar q\,\rho}}&=-\frac{g^{2}}{32\pi^{3}}\int\!d^{3}m\;\ket{\aqk{\alpha}{\lambda}{m}}\,
\frac{t^{b}_{\alpha'\alpha}\rho^{b}(m'-m)}{(p^{+})^{2}}\,\chi^{\dg}_{\lambda'}\Bigg\{
\nonumber\\&\hspace{1.0cm}
\theta(m'^{+}-m^{+})\left[\frac{2}{\bar\Delta-\eps_{p}}
-\frac{p^{+}(\bm\sigma\!\cdot\!\bm m')(\bm\sigma\!\cdot\!\bm m)}{m^{+}m'^{+}(\bar\Delta-\eps_{p})\bar\Delta}\right]
\nonumber\\
&\hspace{1.0cm}-\theta(m^{+}-m'^{+})\left[\frac{2}{\eps_{p}}
-\frac{p^{+}(\bm\sigma\!\cdot\!\bm m')(\bm\sigma\!\cdot\!\bm m)}{m^{+}m'^{+}\eps_{p}\bar\Delta}\right]\Bigg\}\chi_{\lambda},
\label{eq:Pqb-qb}
\end{align}
with $\bar\Delta=\eps_{m'}-\eps_{m}$ and $p=\pm(m'-m)$. We have verified \eqref{eq:Pqb-qb}
explicitly: the cancellation of the instantaneous term proceeds as for the quark, with the
instantaneous matrix element \eqref{eq:MEinstqb} carrying the antiquark charge.

\subsubsection{Summary of results}
\label{sec:sum-q}

The dressed quark state through $\order{g^{2}}$ is
\begin{multline}
\ket{\Psi^{\beta'}_{\lambda'}(l')}=N\ket{\qk{\beta'}{\lambda'}{l'}}
+\ket{(\Psi^{\beta'}_{\lambda'}(l'))_{qg}}+\ket{\qk{\beta'}{\lambda'}{l'}}\otimes\ket{\Psi^{\rm LO}_{g\,\rho}}\\
+\ket{(\Psi^{\beta'}_{\lambda'}(l'))_{q\,\rho}}+\ket{(\Psi^{\beta'}_{\lambda'}(l'))_{qgg}}
+\ket{(\Psi^{\beta'}_{\lambda'}(l'))_{qq\bar q}}+\dots,
\label{eq:sum-q}
\end{multline}
where $\ket{\Psi^{\rm LO}_{g\,\rho}}$ is the leading order Weizs\"acker--Williams component of the
vacuum and the dots denote the remaining disconnected components (the order $g^{2}$ vacuum
components with the quark as a spectator). As in the gluon case, the normalization is fixed by
$\langle\Psi|\Psi\rangle=1$. The $\rho$ dependent overlaps reproduce those of the vacuum, and the
only new contribution is the norm of the $qg$ component. With $z=p^{+}/l'^{+}$ and
$\bm\kappa=\bm p-z\bm l'$ of \eqref{eq:kappaq}, for fixed $\lambda'$,
\begin{equation}
\sum_{\lambda,j}\big|\Ph^{j}_{\lambda\lambda'}(p;l,l')\big|^{2}
=\frac{2\,\bm\kappa^{2}\big[1+(1-z)^{2}\big]}{(l'^{+})^{2}z^{2}(1-z)^{2}},
\end{equation}
and
\begin{multline}
\big\langle(\Psi^{\beta''}_{\lambda''}(l''))_{qg}\big|(\Psi^{\beta'}_{\lambda'}(l'))_{qg}\big\rangle
=\delta^{\beta'\beta''}\delta_{\lambda'\lambda''}\dthree(l'-l'')\,\frac{g^{2}C_{F}}{8\pi^{3}}
\int_{\Lambda/l'^{+}}^{1}\!dz\;\frac{1+(1-z)^{2}}{z}\int\!\frac{d^{2}\kappa}{\bm\kappa^{2}}\\
=\delta\delta\delta\;\frac{g^{2}C_{F}}{8\pi^{3}}\Big[2\log\frac{l'^{+}}{\Lambda}-\frac32\Big]\int\!\frac{d^{2}\kappa}{\bm\kappa^{2}} .
\label{eq:qnorm}
\end{multline}
The longitudinal integral contains the soft gluon logarithm, and the transverse integral is
scaleless and vanishes in dimensional regularization. The normalization of the quark state is
therefore the same $N$ as for all the other sectors, as it must be since $\OM$ is a single
operator.

%----------------------------------------------------------------------
\subsection{Calculation of the quark--antiquark incoming state}
\label{sec:wf-qq}
%----------------------------------------------------------------------

For an incoming pair $\ket{\qk{\beta'}{\lambda_{1}'}{l'}\aqk{\alpha'}{\lambda_{2}'}{m'}}$ the possible
final configurations are the vacuum, one gluon, two gluons, a pair with one gluon and a pair. The
first three test the unitarity relations of the coefficients that annihilate pairs; the last
fixes $\mf D_{3}$.

\subsubsection{Vacuum, one and two gluons outgoing}

At order $g$ the pair annihilates into a gluon:
\begin{multline}
\ket{(\Psi(l',m'))_{g}}=\int\!d^{3}k\;\ket{\gl{a}{i}{k}}\,
\frac{\bra{\gl{a}{i}{k}}H_{gq\bar q}\ket{q\bar q}}{\eps_{l'}+\eps_{m'}-\eps_{k}}\\
=-\frac{g}{8\pi^{3/2}}\int\!d^{3}k\;\ket{\gl{a}{i}{k}}\,
\frac{t^{a}_{\alpha'\beta'}\,\dthree(k-l'-m')\,\big[\Ph^{i}_{\lambda_{1}'\lambda_{2}'}(k;l',m')\big]^{*}}
{\sqrt{k^{+}}\,\big(\eps_{k}-\eps_{l'}-\eps_{m'}\big)} .
\label{eq:Pqq-g}
\end{multline}
Comparing with \eqref{eq:Pg} and \eqref{eq:Omqq} one reads off
$\mf A_{4}=-\mf A_{3}^{\dg}$, which is the first relation of \eqref{eq:U1}; the minus sign is the
sign of the energy denominator. At order $g^{2}$ the pair annihilates into the gluon which is
absorbed by the valence, or annihilates instantaneously on $\rho$. By the reduction formula
\eqref{eq:reduction}, the first is $(-A^{\dg})\mf A_{4}\,\eps_{k}/(\eps_{l'}+\eps_{m'})$ and the second
is the Hermitian conjugate of \eqref{eq:Pvac2} with the opposite sign of the energy
denominator. Their sum reproduces \eqref{eq:UA2} exactly: the instantaneous terms are
anti-Hermitian and cancel between $\mf A_{1}$ and $\mf A_{2}^{\dg}$, while
$-\mf A_{1}^{\dg}+A^{\dg}\mf A_{3}^{\dg}=A^{\dg}\mf A_{3}^{\dg}\,\eps_{k}/(\eps_{l'}+\eps_{m'})+\dots$.
Similarly, the transition $q\bar q\to gg$ computed in LCPT (the time reversed diagrams of
section~\ref{sec:wf-gg}) reproduces the relation \eqref{eq:UB2}, including the terms with $\mf F$
and $\mf G$ which arise from the $t$- and $u$-channel exchanges.

\subsubsection{Quark--antiquark outgoing}

The overall energy denominator is
\begin{equation}
D_{f}\equiv\eps_{l'}+\eps_{m'}-\eps_{l}-\eps_{m}.
\end{equation}
Five connected diagrams contribute at order $g^{2}$.

\paragraph{(a) $s$-channel: the pair annihilates into a gluon which splits again.}
\begin{center}
\begin{tikzpicture}
\draw[qk] (0,1.85) -- (1.4,1.25); \node[left] at (0,1.85) {\scriptsize$\qk{\beta'}{\lambda_{1}'}{l'}$};
\draw[qk] (1.4,1.25) -- (0,0.65); \node[left] at (0,0.65) {\scriptsize$\aqk{\alpha'}{\lambda_{2}'}{m'}$};
\draw[gl] (1.4,1.25) -- (2.5,1.25); \node[above] at (1.95,1.3) {\scriptsize$\gl{a}{i}{k}$};
\draw[qk] (2.5,1.25) -- (3.9,1.85); \node[right] at (3.9,1.85) {\scriptsize$\qk{\beta}{\lambda_{1}}{l}$};
\draw[qk] (3.9,0.65) -- (2.5,1.25); \node[right] at (3.9,0.65) {\scriptsize$\aqk{\alpha}{\lambda_{2}}{m}$};
\draw[cut] (1.95,0.2) -- (1.95,2.25); \draw[cut] (3.2,0.2) -- (3.2,2.25);
\end{tikzpicture}
\end{center}
\begin{equation}
\ket{(\Psi(l',m'))^{1}_{q\bar q}}=\int\! d^{3}k\,d^{3}l\,d^{3}m\;\ket{q\bar q}\;
\frac{\bra{q\bar q}H_{gq\bar q}\ket{\gl{a}{i}{k}}}{D_{f}}\;
\frac{\bra{\gl{a}{i}{k}}H_{gq\bar q}\ket{\qk{\beta'}{\lambda_{1}'}{l'}\aqk{\alpha'}{\lambda_{2}'}{m'}}}
{\eps_{l'}+\eps_{m'}-\eps_{k}} .
\end{equation}
With $(t^{a}_{\beta'\alpha'})^{*}=t^{a}_{\alpha'\beta'}$ and $k=l+m=l'+m'$,
\begin{multline}
\ket{(\Psi(l',m'))^{1}_{q\bar q}}=\frac{g^{2}}{64\pi^{3}}\int\!d^{3}l\,d^{3}m\;\ket{q\bar q}\;
\\
\times t^{a}_{\beta\alpha}t^{a}_{\alpha'\beta'}\;\dthree(l+m-l'-m')\,
\frac{\Ph^{i}_{\lambda_{1}\lambda_{2}}(k;l,m)\,
\big[\Ph^{i}_{\lambda_{1}'\lambda_{2}'}(k;l',m')\big]^{*}}
{k^{+}\,\big(\eps_{l'}+\eps_{m'}-\eps_{k}\big)\,D_{f}} .
\label{eq:D3a}
\end{multline}

\paragraph{(b), (c) $t$-channel gluon exchange.}
\begin{center}
\begin{tikzpicture}
\draw[qk] (0,2.0) -- (1.6,2.0); \node[left] at (0,2.0) {\scriptsize$\qk{\beta'}{\lambda_{1}'}{l'}$};
\draw[qk] (1.6,2.0) -- (3.9,2.0); \node[right] at (3.9,2.0) {\scriptsize$\qk{\beta}{\lambda_{1}}{l}$};
\draw[qk] (3.9,0.3) -- (2.5,0.3); \node[right] at (3.9,0.3) {\scriptsize$\aqk{\alpha}{\lambda_{2}}{m}$};
\draw[qk] (2.5,0.3) -- (0,0.3); \node[left] at (0,0.3) {\scriptsize$\aqk{\alpha'}{\lambda_{2}'}{m'}$};
\draw[gl] (1.6,2.0) -- (2.5,0.3); \node[right] at (2.15,1.2) {\scriptsize$\gl{d}{n}{r}$};
\draw[cut] (2.05,-0.05) -- (2.05,2.35);
\node at (1.95,-0.45) {\scriptsize (b)};
\end{tikzpicture}
\hspace{0.8cm}
\begin{tikzpicture}
\draw[qk] (0,2.0) -- (2.5,2.0); \node[left] at (0,2.0) {\scriptsize$\qk{\beta'}{\lambda_{1}'}{l'}$};
\draw[qk] (2.5,2.0) -- (3.9,2.0); \node[right] at (3.9,2.0) {\scriptsize$\qk{\beta}{\lambda_{1}}{l}$};
\draw[qk] (3.9,0.3) -- (1.6,0.3); \node[right] at (3.9,0.3) {\scriptsize$\aqk{\alpha}{\lambda_{2}}{m}$};
\draw[qk] (1.6,0.3) -- (0,0.3); \node[left] at (0,0.3) {\scriptsize$\aqk{\alpha'}{\lambda_{2}'}{m'}$};
\draw[gl] (1.6,0.3) -- (2.5,2.0); \node[left] at (2.0,1.2) {\scriptsize$\gl{d}{n}{\tilde r}$};
\draw[cut] (2.05,-0.05) -- (2.05,2.35);
\node at (1.95,-0.45) {\scriptsize (c)};
\end{tikzpicture}
\end{center}
In (b) the quark emits the gluon $r=l'-l$ (so that $l'^{+}>l^{+}$) which the antiquark absorbs;
in (c) the antiquark emits $\tilde r=m'-m$ (so that $m'^{+}>m^{+}$) which the quark absorbs.
Using \eqref{eq:MEqg}, \eqref{eq:MEaqg} and their conjugates,
\begin{align}
\ket{(\Psi(l',m'))^{2}_{q\bar q}}&=-\frac{g^{2}}{64\pi^{3}}\int\!d^{3}l\,d^{3}m\;\ket{q\bar q}\;
t^{d}_{\beta\beta'}t^{d}_{\alpha'\alpha}\;\dthree(l+m-l'-m')
\nonumber\\&\hspace{3cm}\times
\frac{\Ph^{n}_{\lambda_{1}\lambda_{1}'}(r;l,l')\,\Ph^{n}_{\lambda_{2}'\lambda_{2}}(r;m',m)}
{r^{+}\,\big(\eps_{l'}-\eps_{l}-\eps_{r}\big)\,D_{f}}\,\theta(r^{+}),
\label{eq:D3b}\\
\ket{(\Psi(l',m'))^{3}_{q\bar q}}&=-\frac{g^{2}}{64\pi^{3}}\int\!d^{3}l\,d^{3}m\;\ket{q\bar q}\;
t^{d}_{\beta\beta'}t^{d}_{\alpha'\alpha}\;\dthree(l+m-l'-m')
\nonumber\\&\hspace{3cm}\times
\frac{\Ph^{n}_{\lambda_{1}\lambda_{1}'}(\tilde r;l,l')\,\Ph^{n}_{\lambda_{2}'\lambda_{2}}(\tilde r;m',m)}
{\tilde r^{+}\,\big(\eps_{m'}-\eps_{m}-\eps_{\tilde r}\big)\,D_{f}}\,\theta(\tilde r^{+}) .
\label{eq:D3c}
\end{align}
The two orderings populate the complementary regions $l'^{+}\gtrless l^{+}$.

\paragraph{(d) Instantaneous gluon exchange between the quark and the antiquark.}
\begin{center}
\begin{tikzpicture}
\draw[qk] (0,2.0) -- (3.9,2.0); \node[left] at (0,2.0) {\scriptsize$\qk{\beta'}{\lambda_{1}'}{l'}$};
\node[right] at (3.9,2.0) {\scriptsize$\qk{\beta}{\lambda_{1}}{l}$};
\draw[qk] (3.9,0.3) -- (0,0.3); \node[left] at (0,0.3) {\scriptsize$\aqk{\alpha'}{\lambda_{2}'}{m'}$};
\node[right] at (3.9,0.3) {\scriptsize$\aqk{\alpha}{\lambda_{2}}{m}$};
\draw[ins] (1.95,2.0) -- (1.95,0.3); \ibar{1.95}{1.15}
\draw[cut] (3.2,-0.05) -- (3.2,2.35);
\end{tikzpicture}
\end{center}
This is the $J_{q}\cdot J_{q}$ term of \eqref{eq:Hqinst}; the exchanged line carries
$P^{+}=l'^{+}-l^{+}$ and the antiquark charge supplies a minus sign:
\begin{equation}
\ket{(\Psi(l',m'))^{4}_{q\bar q}}=-\frac{g^{2}}{(2\pi)^{3}}\int\!d^{3}l\,d^{3}m\;\ket{q\bar q}\;
\frac{t^{a}_{\beta\beta'}t^{a}_{\alpha'\alpha}\,\delta_{\lambda_{1}\lambda_{1}'}\delta_{\lambda_{2}\lambda_{2}'}\,
\dthree(l+m-l'-m')}{(l'^{+}-l^{+})^{2}\;D_{f}} .
\label{eq:D3d}
\end{equation}
The diagrams (b), (c) and (d) together build the full $t$-channel one-gluon exchange between the
quark and the antiquark in light-cone gauge.

\paragraph{(e) Instantaneous $s$-channel gluon exchange.}
\begin{center}
\begin{tikzpicture}
\draw[qk] (0,1.85) -- (1.5,1.25); \node[left] at (0,1.85) {\scriptsize$\qk{\beta'}{\lambda_{1}'}{l'}$};
\draw[qk] (1.5,1.25) -- (0,0.65); \node[left] at (0,0.65) {\scriptsize$\aqk{\alpha'}{\lambda_{2}'}{m'}$};
\draw[ins] (1.5,1.25) -- (2.4,1.25); \ibarh{1.95}{1.25}
\draw[qk] (2.4,1.25) -- (3.9,1.85); \node[right] at (3.9,1.85) {\scriptsize$\qk{\beta}{\lambda_{1}}{l}$};
\draw[qk] (3.9,0.65) -- (2.4,1.25); \node[right] at (3.9,0.65) {\scriptsize$\aqk{\alpha}{\lambda_{2}}{m}$};
\draw[cut] (3.2,0.2) -- (3.2,2.25);
\end{tikzpicture}
\end{center}
The $J_{q}\cdot J_{q}$ interaction also annihilates the incoming pair and creates the outgoing one
instantaneously; the exchanged line carries $P^{+}=k^{+}=l^{+}+m^{+}$:
\begin{equation}
\ket{(\Psi(l',m'))^{5}_{q\bar q}}=\frac{g^{2}}{(2\pi)^{3}}\int\!d^{3}l\,d^{3}m\;\ket{q\bar q}\;
\frac{t^{a}_{\beta\alpha}t^{a}_{\alpha'\beta'}\,\delta_{\lambda_{1}\lambda_{2}}\delta_{\lambda_{1}'\lambda_{2}'}\,
\dthree(l+m-l'-m')}{(k^{+})^{2}\;D_{f}} .
\label{eq:D3e}
\end{equation}
It is the instantaneous part of the $s$-channel exchange (a), in the same way as (d) is that of
(b), (c).

\paragraph{Disconnected contributions.} The exchange of a gluon between the quark (or antiquark)
and the valence charge, with the other member of the pair a spectator, is the coefficient
$\mf P$ ($\bar{\mf P}$) of \eqref{eq:Omqq} computed in section~\ref{sec:wf-q}; it is not part of
$\mf D_{3}$.

\subsubsection{Summary of results}
\label{sec:sum-qq}

The connected $q\bar q\to q\bar q$ amplitude at order $g^{2}$ is the sum of \eqref{eq:D3a}--\eqref{eq:D3e};
its overall denominator $D_{f}$ vanishes on the energy shell and is understood with the
prescription of section~\ref{sec:lcpt}. The remaining components of the dressed pair state are
products of the single quark and single antiquark components and of the order $g$ annihilation
amplitude \eqref{eq:Pqq-g}.

%----------------------------------------------------------------------
\subsection{Calculation of the quark--gluon incoming state}
\label{sec:wf-qg}
%----------------------------------------------------------------------

Finally, for an incoming state $\ket{\qk{\beta'}{\lambda'}{l'}\gl{c}{k}{q}}$ the connected order
$g^{2}$ transition into $\ket{\qk{\beta}{\lambda}{l}\gl{d}{n}{r}}$ is soft quark--gluon (Compton)
scattering. The light-cone time ordered contributions are: the $s$-channel (the quark absorbs
the gluon and re-emits one); the $u$-channel in its two orderings (the quark first emits the
gluon $r$ and then absorbs $q$, or the gluon $q$ first splits into a pair whose antiquark
annihilates with the incoming quark into the gluon $r$); the $t$-channel in its two orderings (the
gluon $q$ splits into two gluons one of which is absorbed by the quark, or the quark emits a gluon
which merges with $q$); the instantaneous quark exchange; and the instantaneous $J_{q}\cdot J_{g}$
gluon exchange. In addition there are the $\rho$ dependent products in which the incoming gluon is
absorbed by the valence while the quark emits the outgoing gluon, or the valence emits the outgoing
gluon while the quark absorbs the incoming one; by \eqref{eq:factor} these factorize. With the
reduction formula \eqref{eq:reduction}, each contribution is the product of two of the first order
amplitudes already computed, weighted by the corresponding ratio of energy denominators, or an
instantaneous matrix element divided by $D_{f}=\eps_{l'}+\eps_{q}-\eps_{l}-\eps_{r}$. The resulting
coefficient $\mf Q_{2}$ is given in section~\ref{sec:coeff}; like $\mf D_{3}$, its $\rho$
independent part has an on-shell pole.
